% !TEX root = ../main.tex
\chapter{TINJAUAN PUSTAKA}
\label{bab:tinjauan}
% Pedoman 9.2: teori, kajian literatur, kesenjangan, kerangka pemikiran,
% definisi operasional. Bab ini harus berakhir pada celah yang diklaim di Bab 1.

\section{Landasan Teori}
\label{sec:teori}
\subsection{Manajemen Kerentanan dan CVE}
...
\subsection{\asing{Large Language Model} dan Kuantisasi}
...
\subsection{\asing{Serving} LLM dengan vLLM}
...
\subsection{Sistem Multi-Agen dan Orkestrasi dengan n8n}
...
\subsection{Penggabungan Jawaban Beberapa LLM}
% Self-consistency, majority/weighted voting, multi-agent debate,
% LLM-as-a-judge/chairman. Bagian teori inti: beri ruang paling banyak.
...
\subsection{Evaluasi Klasifikasi}
% Confusion matrix, precision, recall, F1, Fleiss' kappa, uji McNemar.
...

\section{Penelitian Terkait}
\label{sec:penelitian-terkait}
% Kelompokkan per tema, bukan satu paragraf per paper.
\subsection{\asing{Multi-Agent Debate} dan \asing{LLM Council}}
...
\subsection{\asing{Ensemble} dan \asing{Voting} pada LLM}
...
\subsection{\asing{LLM-as-a-Judge}}
...
\subsection{LLM dan Agen untuk Keamanan Siber}
...
\subsection{Model Kecil dan \asing{Deployment} Lokal}
...

% Tabel perbandingan penelitian terkait (contoh format tabel).
\begin{table}[htbp]
  \caption{Perbandingan penelitian terkait}
  \label{tab:perbandingan}
  \centering\small\setstretch{1}
  \begin{tabularx}{\textwidth}{@{}l X l l X@{}}
    \toprule
    Penelitian & Domain & Jumlah model & Lokal/\asing{cloud} & Hasil utama \\
    \midrule
    ... & ... & ... & ... & ... \\
    Penelitian ini & Verifikasi kerentanan & 1 vs 3 & Lokal & ... \\
    \bottomrule
  \end{tabularx}
\end{table}

\section{Analisis Kesenjangan Penelitian}
\label{sec:kesenjangan}
...

\section{Kerangka Pemikiran}
\label{sec:kerangka}
% Gambar model konseptual: masalah -> intervensi -> hasil yang diukur.
\begin{figure}[htbp]
  \centering
  \fbox{\parbox[c][4cm][c]{0.8\textwidth}{\centering Gambar kerangka pemikiran}}
  \caption{Kerangka pemikiran penelitian}
  \label{fig:kerangka}
\end{figure}

\section{Definisi Operasional}
\label{sec:definisi-operasional}
\begin{table}[htbp]
  \caption{Definisi operasional}
  \label{tab:definisi-operasional}
  \centering\small\setstretch{1}
  \begin{tabularx}{\textwidth}{@{}l X X l@{}}
    \toprule
    Istilah & Definisi & Cara mengukur & Satuan \\
    \midrule
    \asing{True positive} & ... & ... & kasus \\
    \asing{False positive} & ... & ... & kasus \\
    F1 & ... & ... & -- \\
    Latensi p50/p95 & ... & ... & detik \\
    Total token & ... & ... & token \\
    VRAM puncak & ... & ... & GB \\
    Fleiss' kappa & ... & ... & -- \\
    \bottomrule
  \end{tabularx}
\end{table}
